\section{Introduction}
Mixture-of-experts (MoE) models activate only a small part of their parameters for each token \citep{switch,deepseekv3}. DeepSeek-V4.1-Flash has 754B parameters in 40 layers, and each layer has 384 routed experts, of which a token uses six \citep{deepseek41flash}. In the data center, DeepSeek serves this model with expert parallelism (EP) \citep{deepinfer}. The experts of a layer are spread over many GPUs, and after routing each token is sent to the GPUs that hold its experts. The cost of this communication is hidden by overlapping two micro-batches, and frequently used experts are replicated.

This paper asks whether the same model can run on a MacBook Air with 24\,GB of memory, and which of these techniques still help there. The converted model takes 338\,GB, so it cannot be held in memory. However, the work for one token is much smaller. The dense layers take 7.6\,GB at 8 bits per weight, and the six experts per layer that a token needs take about 3.3\,GB in total. The laptop also has a second device that can store experts: an SSD that reads about 3\,GB/s. We therefore treat the SSD as the expert-holding device, the role that other GPUs play under EP.

We built \method{} around this idea (Figure~\ref{fig:overview}). After routing, the engine reads each missing expert from the SSD with one \texttt{preadv} call directly into a slot of a GPU buffer, without copies and without the OS page cache, and runs all experts of the layer in one gathered quantized matmul. Recently used experts stay in memory as a cache, which plays the role of expert replication.

Measuring this system showed where the time goes. A decoding token forms a strict chain in every layer: the router decides which experts to read only after attention, and the next layer can start only after the experts are computed. The SSD is therefore idle while the GPU works, and the GPU waits while the SSD reads. The cluster's micro-batch overlap helps little here: computing the cached experts while the missing ones load made decode 5\% faster. Prefetching the next layer's experts \citep{moeoffload} should fill the idle time, but in our measurements it made decode 5--11\% slower. An interleaved ablation shows why: the guesses compete with reads the token is waiting for. \emph{Idle-window speculation} issues the guesses only after the current layer's misses have arrived, so that they use the SSD while the GPU computes; it also inserts them at the cold end of the cache, which adds a small gain within our confidence intervals. This makes decode 11--27\% faster at every cache size we tested and leaves the output unchanged. In three traced tokens per setting, the SSD is busy 62--65\% of a decode step without speculation and 73--89\% with it.

The engine follows the official reference implementation op by op. It includes compressed sparse attention, Engram memory, Hyper-Connections \citep{hyperconnections}, the vision encoder, DSpark drafting and the FP8/FP4 roundings used in training, and our tests run the official PyTorch model as an oracle.

In summary, we make the following contributions:
\begin{itemize}
  \item \textbf{Expert parallelism on one laptop.} We map EP to one GPU and one SSD: experts move to the GPU through zero-copy reads into fixed buffer slots, and an LRU cache replaces expert replication. We are not aware of another system that runs every component of DeepSeek-V4.1-Flash, including Engram, Hyper-Connections and vision, within 24\,GB with parity tests against the reference (Sections~\ref{sec:design} and~\ref{sec:impl}).
  \item \textbf{A cost model and a study of parallelism.} We model a decode step as a chain of compute and SSD reads. A two-parameter fit matches measured decode speed within 2\% across cache sizes; applied to the speculation runs it shows that speculative reads are not free, and traces show that they mostly fill time when the SSD was idle. We test parallelism at six levels and find which kinds help: those that give the SSD work while the chain leaves it idle, without taking memory or bandwidth from demand reads (Sections~\ref{sec:chain} and~\ref{sec:parallel-eval}).
  \item \textbf{Idle-window speculation.} We find why next-layer expert prefetching \citep{moeoffload} slows decode on a memory-bound laptop: the guesses delay reads the token is waiting for. Issuing guesses only after the demand reads have landed, a known prefetch-control idea, turns the slowdown into a speed-up; an ablation separates it from cold insertion, which adds little (Section~\ref{sec:spec} and~\ref{sec:parallel-eval}).
  \item \textbf{An exactly equal, faster indexer.} We prove that, when the scales of each FP4 block product differ by at most $2^9$ the products stay in the normal FP32 range and the matrix units accumulate in FP32, the FP32 indexer score of DSA-style indexers can be computed with BF16 matrix units without changing a bit. Together with a fused head reduction, this makes the indexer about $2\times$ faster in long prefill ($1.8$--$2.1\times$ on three layers, $3.5\times$ on one). We also measure DoSA's exact screening \citep{dosa} on this model's compressed two-level indexer, where it saves little at 16--32K keys (Section~\ref{sec:indexer}).
\end{itemize}
On a MacBook Air (M4, 24\,GB), \method{} decodes 1.23 tokens/s with a 5.5\,GB expert cache, which keeps 9\,GB free for macOS, and 1.35 tokens/s with 6.5\,GB. Its next-token predictions agree with DeepSeek's API on 86.0\% of open-ended continuation tokens, and it answers a retrieval question correctly at every tested length up to 16K tokens. The code, converter, tests and measurement scripts are open source.

\begin{figure}[t]
\centering
\resizebox{\linewidth}{!}{%
\begin{tikzpicture}[
  font=\small,
  blk/.style={draw=black!45, rounded corners=3pt, fill=white, align=center, minimum height=1.25cm, inner sep=4pt},
  arr/.style={-{Stealth[length=6pt]}, draw=dsgray, line width=0.8pt},
  rd/.style={-{Stealth[length=6pt]}, draw=dsorange, line width=1pt},
  region/.style={draw, dashed, rounded corners=6pt, inner sep=8pt}]
  \node[blk, text width=3.7cm] (att) {\textbf{Attention}\\window + compressed sparse,\\indexer, mHC, Engram};
  \node[blk, text width=2.7cm, right=0.6cm of att] (rt) {\textbf{Router}\\6 of 384 experts\\+ shared expert};
  \node[blk, text width=3.1cm, right=0.6cm of rt] (cache) {\textbf{Expert cache}\\hot experts replicated\\in GPU memory};
  \node[blk, text width=3.1cm, below=0.55cm of cache] (gq) {\textbf{gather\_qmm}\\all experts of the layer,\\one quantized matmul};
  \node[blk, text width=2.9cm, right=2.7cm of cache] (exp) {\textbf{Routed experts}\\$40\times384$, 2/3/4-bit\\208\,GB};
  \node[blk, text width=2.9cm, below=0.55cm of exp] (eg) {\textbf{Engram tables}\\$2\times384$M rows, 4-bit\\114\,GB};
  \coordinate (low) at ($(gq.south)+(0,-0.55)$);
  \draw[arr] (att) -- (rt);
  \draw[arr] (rt) -- node[above, font=\footnotesize, text=black!70] {ids} (cache);
  \draw[arr] (cache) -- (gq);
  \draw[rd] (exp) -- node[above, text=dsorange, font=\footnotesize\bfseries] {miss} node[below, align=center, font=\footnotesize, text=black!70] {one \texttt{preadv},\\zero-copy} (cache);
  \draw[rd, dashed] (eg.south) |- (low) -| node[pos=0.25, above, font=\footnotesize, text=black!70] {Engram rows (layers 1, 14)} (att.south);
  \begin{scope}[on background layer]
    \node[region, draw=dsblue, fill=dsblue!6, fit=(att)(rt)(cache)(gq)(low), label={[anchor=south west, text=dsblue, font=\footnotesize\bfseries]north west:GPU AND UNIFIED MEMORY (24\,GB): 7.6\,GB WEIGHTS + EXPERT CACHE}] {};
    \node[region, draw=dsorange, fill=dsorange!6, fit=(exp)(eg), label={[anchor=south west, text=dsorange, font=\footnotesize\bfseries]north west:SSD: THE EXPERT PEER}] {};
  \end{scope}
\end{tikzpicture}}
\caption{Overview of \method{}. Dense weights and recently used experts are kept in unified memory. After routing, missing experts are read from the SSD directly into slots of preallocated GPU buffers. Within a layer, attention, router, reads and expert computation run one after another.}
\label{fig:overview}
\end{figure}
